\usepackage{mathtools}
\usepackage{stmaryrd}
\usepackage{textcomp}
% type setting
% -----------------------------------------------------------------------------
% \usepackage[ngerman]{babel}

% fonts
% -----------------------------------------------------------------------------
\usefonttheme{professionalfonts}

% slide title and horizontal line
% -----------------------------------------------------------------------------
% Feste Titelhöhe ohne negative Verschiebung außerhalb der Folie.
\setbeamertemplate{headline}{}
\setbeamertemplate{frametitle}{%
  \vspace{2mm}%
  \begin{beamercolorbox}[wd=\textwidth,ht=5mm,dp=1mm]{frametitle}%
    \color{black}\large\insertframetitle
  \end{beamercolorbox}%
  \vspace{-1mm}\par
  \color{black}\rule{\textwidth}{0.4pt}\par
  \vspace{2mm}%
}

% slide number
% -----------------------------------------------------------------------------
\setbeamertemplate{navigation symbols}{}
\setbeamertemplate{footline}
{
\vskip5pt
\vskip2pt
\makebox[123mm]{\hspace{7.5mm}
\hfill Grundlagen der Testtheorie $\vert$
\copyright $ $ 2027 Dirk Ostwald CC BY 4.0 $\vert$
Folie \insertframenumber}
\vskip4pt
}

% block color scheme
% -----------------------------------------------------------------------------
% colors
\definecolor{white}{RGB}{255,255,255}
\definecolor{grey}{RGB}{235,235,235}
\definecolor{lightgrey}{RGB}{245,245,245}
\definecolor{LightBlue}{RGB}{220,220,255}
\definecolor{darkblue}{RGB}{51, 51, 153}

% definitions and theorems
\setbeamercolor{block title}{fg = black, bg = grey}
\setbeamercolor{block body}{fg = black, bg = lightgrey}

% more colors
\definecolor{green}{HTML}{209C05}
\definecolor{lightgreen}{HTML}{85E62C}
\definecolor{yellow}{HTML}{EBFF0A}
\definecolor{orange}{HTML}{EE6F27}
\definecolor{red}{HTML}{FF0A0A}

% general line spacing
% -----------------------------------------------------------------------------
\linespread{1.3}

% local line spacing
% -----------------------------------------------------------------------------
\usepackage{setspace}

% colors
% -----------------------------------------------------------------------------
\usepackage{color}
\usepackage{xcolor}

% justified text
% -----------------------------------------------------------------------------
\usepackage{ragged2e}
\usepackage{etoolbox}
\apptocmd{\frame}{}{\justifying}{}
\usepackage{ragged2e}
\renewcommand{\raggedright}{\justifying}

% bullet point lists
% -----------------------------------------------------------------------------
\setbeamertemplate{itemize item}[circle]
\setbeamertemplate{itemize subitem}[circle]
\setbeamertemplate{itemize subsubitem}[circle]
\setbeamercolor{itemize item}{fg = black}
\setbeamercolor{itemize subitem}{fg = black}
\setbeamercolor{itemize subsubitem}{fg = black}
\setbeamercolor{enumerate item}{fg = black}
\setbeamercolor{enumerate subitem}{fg = black}
\setbeamercolor{enumerate subsubitem}{fg = black}
\setbeamerfont{itemize/enumerate body}{}
\setbeamerfont{itemize/enumerate subbody}{size = \normalsize}
\setbeamerfont{itemize/enumerate subsubbody}{size = \normalsize}

% tables
% -----------------------------------------------------------------------------
\usepackage{tabularx}

% color links
% -----------------------------------------------------------------------------
\usepackage{hyperref}
\definecolor{urls}{RGB}{204,0,0}
\hypersetup{colorlinks = true, citecolor = darkblue, urlcolor = darkblue, linkcolor = darkblue}


% additional math commands
% -----------------------------------------------------------------------------
\usepackage{amsmath}
\usepackage{bm}
\newcommand{\niton}{\not\owns}
\newcommand{\ups}{\upsilon}
\DeclareMathOperator*{\intinf}{\int_{-\infty}^{\infty}}

% text highlighting
% -----------------------------------------------------------------------------
\usepackage{soul}
\makeatletter
\let\HL\hl
\renewcommand\hl{%
  \let\set@color\beamerorig@set@color
  \let\reset@color\beamerorig@reset@color
  \HL}
\makeatother

% equation highlighting
% -----------------------------------------------------------------------------
\newcommand{\highlight}[2][yellow]{\mathchoice%
  {\colorbox{#1}{$\displaystyle#2$}}%
  {\colorbox{#1}{$\textstyle#2$}}%
  {\colorbox{#1}{$\scriptstyle#2$}}%
  {\colorbox{#1}{$\scriptscriptstyle#2$}}}%

% additional mathematical operators
% ------------------------------------------------------------------------------
\DeclareMathOperator*{\argmax}{arg\,max}
\DeclareMathOperator*{\argmin}{arg\,min}


% spacing between bibliography entries
% -----------------------------------------------------------------------------
\AtBeginDocument{%
  \apptocmd{\CSLReferences}{\setlength{\itemsep}{1pt}}{}{}%
}

% Mathematical notation and quotation marks used in the course
\providecommand{\glqq}{\quotedblbase}
\providecommand{\grqq}{\textquotedblleft}
\providecommand{\glq}{\quotesinglbase}
\providecommand{\grq}{\textquoteleft}
\newcommand*{\thmref}[1]{\hyperref[{#1}]{Theorem \textcolor{darkblue}{\nameref*{#1}}}}
\newcommand*{\defref}[1]{\hyperref[{#1}]{Definition \textcolor{darkblue}{\nameref*{#1}}}}
\newcommand{\mrm}{\mathrm}
\newcommand{\mbb}{\mathbb}
\newcommand{\eps}{\varepsilon}
\DeclarePairedDelimiter{\ib}{\llbracket}{\rrbracket}
\DeclareMathOperator{\tr}{\mbox{tr}}

% Links to the other slide decks
\hypersetup{filecolor=darkblue, runcolor=darkblue}
